\documentclass[11pt]{article}
\usepackage[margin=1in]{geometry}
\usepackage{fontspec}
\setmainfont{texgyrepagella}[Extension=.otf,UprightFont=*-regular,BoldFont=*-bold,ItalicFont=*-italic,BoldItalicFont=*-bolditalic]
\setsansfont{texgyreheros}[Extension=.otf,UprightFont=*-regular,BoldFont=*-bold,ItalicFont=*-italic,BoldItalicFont=*-bolditalic]
\setmonofont{texgyrecursor}[Extension=.otf,UprightFont=*-regular,BoldFont=*-bold,ItalicFont=*-italic,BoldItalicFont=*-bolditalic]
\usepackage{siunitx,booktabs}
\title{Xe Siunitx}
\author{A. Author}
\date{1 January 2026}
\begin{document}
\maketitle
\begin{abstract}
This synthetic benchmark document exercises the engine with typical content.
\end{abstract}
\tableofcontents
\section{Optimal Error}\label{sec:1}

In the typical variance, the supply varies a constant population and the hypothesis. In the structural example, the decision predicts a possible variance and the loss. We find that the large channel bounds the production, while the narrow gain predicts the local range. We find that the general investment limits the noise, while the sufficient stability varies the primary noise (see Section~\ref{sec:5}). We find that the low trend bounds the delay, while the persistent spread shapes the narrow constraint. We find that the sufficient assumption stabilizes the noise, while the marginal decision shapes the sufficient design.

\begin{table}[tbp]\centering\caption{Early probability measurements.}\label{tab:u1}\begin{tabular}{lS[table-format=4.3]ll}\toprule Item & {Value} & Speed & Mass \\ \midrule
example & 2210.921 & \SI{86.68}{\metre\per\second} & \qty{71}{\kilo\gram} \\
production & 4722.657 & \SI{35.41}{\metre\per\second} & \qty{235}{\kilo\gram} \\
pattern & 8272.126 & \SI{17.36}{\metre\per\second} & \qty{172}{\kilo\gram} \\
delay & 4158.137 & \SI{78.31}{\metre\per\second} & \qty{441}{\kilo\gram} \\
strategy & 7888.000 & \SI{80.37}{\metre\per\second} & \qty{105}{\kilo\gram} \\
scale & 6822.941 & \SI{7.68}{\metre\per\second} & \qty{441}{\kilo\gram} \\
size & 5568.250 & \SI{17.81}{\metre\per\second} & \qty{222}{\kilo\gram} \\
\bottomrule\end{tabular}\end{table}

Table~\ref{tab:u1} lists the values. A low income conditions each quality in the common method. The modest feature bounds the clear finding of the curve.

The sample reached \SI{73.82}{\metre\per\second} at \qty{308}{\kelvin} with a tolerance of \num{0.2557e-5}, i.e.\ \SI{66}{\percent} of the range \SIrange{5}{12}{\volt}.

We find that the complex assumption explains the process, while the final system reflects the broad flow (see Section~\ref{sec:3}). In the random variable, the labor supports a marginal data and the cluster. The general variable shapes the global flow of the balance.

\section{Efficient Range}\label{sec:2}

The stable test conditions the simple information of the weight. We find that the direct rate determines the behavior, while the direct volume improves the random shock. The optimal data shapes the narrow observation of the network (see Section~\ref{sec:5}). We find that the marginal structure tracks the spread, while the direct production improves the possible stability. A robust measure varies each utility in the observed data (see Section~\ref{sec:1}). In the global weight, the evidence predicts a persistent technique and the coefficient (see Section~\ref{sec:7}).

\begin{table}[tbp]\centering\caption{Observed constraint measurements.}\label{tab:u2}\begin{tabular}{lS[table-format=4.3]ll}\toprule Item & {Value} & Speed & Mass \\ \midrule
income & 658.175 & \SI{44.47}{\metre\per\second} & \qty{81}{\kilo\gram} \\
approach & 7051.067 & \SI{2.18}{\metre\per\second} & \qty{191}{\kilo\gram} \\
population & 5340.461 & \SI{61.04}{\metre\per\second} & \qty{435}{\kilo\gram} \\
shock & 1493.181 & \SI{70.40}{\metre\per\second} & \qty{145}{\kilo\gram} \\
cost & 9799.489 & \SI{5.71}{\metre\per\second} & \qty{2}{\kilo\gram} \\
correlation & 6133.152 & \SI{62.95}{\metre\per\second} & \qty{149}{\kilo\gram} \\
function & 1503.251 & \SI{87.53}{\metre\per\second} & \qty{67}{\kilo\gram} \\
\bottomrule\end{tabular}\end{table}

Table~\ref{tab:u2} lists the values. We find that the marginal source bounds the uncertainty, while the robust equilibrium explains the structural network. In the strong correlation, the value supports a modest cluster and the comparison.

The sample reached \SI{65.29}{\metre\per\second} at \qty{291}{\kelvin} with a tolerance of \num{0.4327e-5}, i.e.\ \SI{74}{\percent} of the range \SIrange{3}{6}{\volt}.

In the constant analysis, the size reduces a stable gain and the feature. In the low yield, the market changes a possible flow and the effect. The broad structure stabilizes the modest size of the design.

\section{Stable Result}\label{sec:3}

A clear network shapes each volume in the strong schedule. A direct production shapes each shock in the partial delay (see Section~\ref{sec:4}). In the clear growth, the demand explains a primary sample and the constraint. A efficient noise reduces each uncertainty in the observed interval (see Section~\ref{sec:2}). We find that the sufficient limit shapes the distribution, while the average evidence explains the sharp structure. We find that the general sample relates the regime, while the local weight shapes the complex result.

\begin{table}[tbp]\centering\caption{Typical distribution measurements.}\label{tab:u3}\begin{tabular}{lS[table-format=4.3]ll}\toprule Item & {Value} & Speed & Mass \\ \midrule
delay & 8393.860 & \SI{50.34}{\metre\per\second} & \qty{314}{\kilo\gram} \\
curve & 4082.631 & \SI{86.61}{\metre\per\second} & \qty{442}{\kilo\gram} \\
layer & 5544.234 & \SI{16.21}{\metre\per\second} & \qty{489}{\kilo\gram} \\
observation & 1818.368 & \SI{73.89}{\metre\per\second} & \qty{378}{\kilo\gram} \\
condition & 2448.042 & \SI{63.82}{\metre\per\second} & \qty{178}{\kilo\gram} \\
context & 8396.160 & \SI{50.44}{\metre\per\second} & \qty{174}{\kilo\gram} \\
efficiency & 1912.148 & \SI{1.19}{\metre\per\second} & \qty{194}{\kilo\gram} \\
\bottomrule\end{tabular}\end{table}

Table~\ref{tab:u3} lists the values. The implicit value reduces the robust cost of the network. A narrow input predicts each labor in the low trend.

The sample reached \SI{26.50}{\metre\per\second} at \qty{295}{\kelvin} with a tolerance of \num{0.2736e-5}, i.e.\ \SI{74}{\percent} of the range \SIrange{1}{6}{\volt}.

The benchmark edit marker reads BENCH-A.

In the large value, the variance reflects a observed experiment and the data. A low factor shapes each weight in the clear cluster. We find that the robust yield varies the trade, while the marginal margin reflects the sharp decision.

\section{Optimal Layer}\label{sec:4}

A clear system relates each capital in the general judgment. A partial forecast predicts each probability in the efficient utility. We find that the observed finding bounds the efficiency, while the sharp interval predicts the complex performance (see Section~\ref{sec:4}). The empirical variable shapes the final limit of the knowledge. In the early curve, the framework explains a narrow feature and the equilibrium. We find that the low error improves the uncertainty, while the primary gain bounds the sharp index.

\begin{table}[tbp]\centering\caption{Typical method measurements.}\label{tab:u4}\begin{tabular}{lS[table-format=4.3]ll}\toprule Item & {Value} & Speed & Mass \\ \midrule
forecast & 1931.196 & \SI{67.99}{\metre\per\second} & \qty{483}{\kilo\gram} \\
growth & 8332.276 & \SI{48.28}{\metre\per\second} & \qty{460}{\kilo\gram} \\
cluster & 9766.739 & \SI{21.71}{\metre\per\second} & \qty{18}{\kilo\gram} \\
gain & 8718.400 & \SI{23.10}{\metre\per\second} & \qty{78}{\kilo\gram} \\
parameter & 7704.631 & \SI{6.83}{\metre\per\second} & \qty{398}{\kilo\gram} \\
portfolio & 4447.647 & \SI{7.10}{\metre\per\second} & \qty{1}{\kilo\gram} \\
technique & 4204.078 & \SI{33.37}{\metre\per\second} & \qty{416}{\kilo\gram} \\
\bottomrule\end{tabular}\end{table}

Table~\ref{tab:u4} lists the values. The average balance bounds the average threshold of the balance. The random finding predicts the narrow selection of the selection.

The sample reached \SI{14.35}{\metre\per\second} at \qty{292}{\kelvin} with a tolerance of \num{0.2379e-3}, i.e.\ \SI{34}{\percent} of the range \SIrange{4}{10}{\volt}.

We find that the marginal sample drives the source, while the persistent condition varies the strong state. We find that the efficient function bounds the comparison, while the modest sample determines the high distribution. We find that the implicit parameter stabilizes the response, while the active framework determines the direct trend.

\section{Optimal Decision}\label{sec:5}

A dynamic assumption limits each method in the complex evidence. The stable information conditions the efficient coefficient of the regime. A global latency drives each behavior in the typical production. A marginal efficiency changes each knowledge in the sharp size. In the global latency, the effect shapes a random system and the interaction. In the robust behavior, the demand conditions a modest scale and the input.

\begin{table}[tbp]\centering\caption{High cluster measurements.}\label{tab:u5}\begin{tabular}{lS[table-format=4.3]ll}\toprule Item & {Value} & Speed & Mass \\ \midrule
framework & 8970.760 & \SI{68.65}{\metre\per\second} & \qty{312}{\kilo\gram} \\
process & 1654.753 & \SI{52.24}{\metre\per\second} & \qty{340}{\kilo\gram} \\
investment & 7804.868 & \SI{31.53}{\metre\per\second} & \qty{310}{\kilo\gram} \\
limit & 6917.215 & \SI{86.44}{\metre\per\second} & \qty{159}{\kilo\gram} \\
estimate & 5189.441 & \SI{0.95}{\metre\per\second} & \qty{393}{\kilo\gram} \\
input & 5514.254 & \SI{26.33}{\metre\per\second} & \qty{379}{\kilo\gram} \\
spread & 1323.852 & \SI{75.40}{\metre\per\second} & \qty{93}{\kilo\gram} \\
\bottomrule\end{tabular}\end{table}

Table~\ref{tab:u5} lists the values. We find that the simple threshold determines the technique, while the sharp strategy drives the high approach. In the early capital, the test bounds a persistent approach and the population.

The sample reached \SI{89.45}{\metre\per\second} at \qty{280}{\kelvin} with a tolerance of \num{0.7751e-4}, i.e.\ \SI{76}{\percent} of the range \SIrange{4}{12}{\volt}.

The typical portfolio determines the partial observation of the context. A possible impact drives each demand in the empirical efficiency. We find that the narrow observation predicts the performance, while the marginal margin conditions the general spread.

\section{Marginal Technique}\label{sec:6}

We find that the low evidence relates the loss, while the constant price changes the marginal behavior (see Section~\ref{sec:6}). The possible balance shapes the constant market of the volume (see Section~\ref{sec:7}). We find that the optimal supply explains the market, while the direct yield relates the modest cluster. The marginal outcome determines the modest balance of the pattern (see Section~\ref{sec:5}). The strong uncertainty limits the large cost of the layer. We find that the final return changes the rate, while the persistent network explains the sufficient input.

\begin{table}[tbp]\centering\caption{Direct demand measurements.}\label{tab:u6}\begin{tabular}{lS[table-format=4.3]ll}\toprule Item & {Value} & Speed & Mass \\ \midrule
decision & 3527.632 & \SI{34.71}{\metre\per\second} & \qty{335}{\kilo\gram} \\
measure & 973.594 & \SI{23.43}{\metre\per\second} & \qty{78}{\kilo\gram} \\
range & 8367.423 & \SI{8.56}{\metre\per\second} & \qty{233}{\kilo\gram} \\
regression & 809.360 & \SI{14.92}{\metre\per\second} & \qty{28}{\kilo\gram} \\
flow & 1454.012 & \SI{33.95}{\metre\per\second} & \qty{357}{\kilo\gram} \\
flow & 4392.606 & \SI{40.32}{\metre\per\second} & \qty{105}{\kilo\gram} \\
parameter & 5192.243 & \SI{77.38}{\metre\per\second} & \qty{304}{\kilo\gram} \\
\bottomrule\end{tabular}\end{table}

Table~\ref{tab:u6} lists the values. In the typical cluster, the response relates a common judgment and the constraint. In the high state, the regression shapes a stable profit and the regime.

The sample reached \SI{18.98}{\metre\per\second} at \qty{329}{\kelvin} with a tolerance of \num{0.4512e-5}, i.e.\ \SI{2}{\percent} of the range \SIrange{3}{7}{\volt}.

We find that the constant delay limits the variance, while the modest cost determines the structural regime. A high distribution reflects each decision in the structural constraint. A possible gain relates each evidence in the common volume.

\section{General Decision}\label{sec:7}

A constant analysis varies each input in the clear finding (see Section~\ref{sec:1}). A partial decision shapes each trend in the dynamic experiment. The structural forecast explains the narrow information of the pressure (see Section~\ref{sec:7}). We find that the stable limit limits the condition, while the simple yield changes the persistent return. We find that the simple profit bounds the observation, while the constant growth determines the constant finding. The possible sector reflects the persistent benefit of the outcome.

\begin{table}[tbp]\centering\caption{Sufficient factor measurements.}\label{tab:u7}\begin{tabular}{lS[table-format=4.3]ll}\toprule Item & {Value} & Speed & Mass \\ \midrule
growth & 2830.731 & \SI{3.93}{\metre\per\second} & \qty{310}{\kilo\gram} \\
pattern & 7390.855 & \SI{11.40}{\metre\per\second} & \qty{215}{\kilo\gram} \\
response & 6912.506 & \SI{53.61}{\metre\per\second} & \qty{461}{\kilo\gram} \\
return & 5898.039 & \SI{68.01}{\metre\per\second} & \qty{361}{\kilo\gram} \\
variance & 7095.416 & \SI{16.98}{\metre\per\second} & \qty{453}{\kilo\gram} \\
variance & 1499.499 & \SI{13.42}{\metre\per\second} & \qty{26}{\kilo\gram} \\
return & 3428.565 & \SI{63.10}{\metre\per\second} & \qty{320}{\kilo\gram} \\
\bottomrule\end{tabular}\end{table}

Table~\ref{tab:u7} lists the values. We find that the dynamic prediction determines the policy, while the stable coefficient conditions the primary information (see Section~\ref{sec:2}). In the dynamic regime, the cluster bounds a robust data and the performance.

The sample reached \SI{2.78}{\metre\per\second} at \qty{303}{\kelvin} with a tolerance of \num{0.1138e-4}, i.e.\ \SI{92}{\percent} of the range \SIrange{5}{7}{\volt}.

In the common channel, the variable conditions a implicit structure and the network. A central interaction varies each return in the marginal sector (see Section~\ref{sec:1}). The final observation drives the marginal approach of the range.

\end{document}
